\section{A realization statement and its mirror test objects}
\label{sec:setup}

The first step is to realize a rational middle-dimensional class by a
perfect complex whose second self-extension group is small.  The
realization will preserve its pairings with a fixed collection of
Lagrangian test objects.  We specify those objects here and then use the
realization statement to construct the algebraic tensors.  Its proof occupies
Sections~\ref{sec:topology}--\ref{sec:realization} in the second part
of the paper.

\subsection{The symplectic spaces and the realization statement}

Let $Y$ be either a smooth quartic K3 surface with the integral
symplectic form $\omega_Y$ representing its hyperplane class, or a point.
In the latter case all $Y$-factors below are omitted.  Set
\begin{equation}
 T_g=\R^{2g}/\Z^{2g},\qquad
 \omega_T=\sum_{j=1}^g q_j\wedge p_j,\qquad
 D=Y\times T_g,\qquad 2n=\dim_{\R}D.
 \label{setup:spaces}
\end{equation}
Here $q_j,p_j$ denote the coordinate one-forms; we use the same letters
for coordinates when describing graphs.  Write $\omega$ for the product
symplectic form.  Throughout this construction $g$ is even and $n\ge12$.
All branes have integer Maslov gradings, spin structures with zero
background class, and the orientations determined by their gradings.

The ordinary mirrors are the mirror quartic $B$ of $Y$ and the product
$E^g$, where $E$ is the Tate elliptic curve; set $B$ equal to a point
when $Y$ is a point.  These varieties are initially
defined over Novikov fields.  The change of parameter in quartic
homological mirror symmetry is included in the definition of $B$.
We allow further field extensions and then an abstract field embedding
into $\C$.  Thus passage to a complex variety means extension of scalars
along a field embedding, not evaluation of a Novikov series at a complex
number.  All fields used below have cardinality at most the continuum.

Choose a finite collection of graded Lagrangian spheres $\Sigma_\ell\subset Y$
which split-generates the quartic Fukaya category and whose classes span
\begin{equation}
 P_Y=\omega_Y^\perp\subset H^2(Y,\Q).
 \label{setup:primitive-space}
\end{equation}
Their existence is recalled below.  For the point use one test object,
with the unit convention.  In $T_g$ choose finitely many graphs
\begin{equation}
 G_M=\{p=Mq\},\qquad M\in\operatorname{Sym}_g(\Z),
 \label{setup:graphs}
\end{equation}
and make small Hamiltonian perturbations when transversality is required.
The graph gradings vary continuously with the symmetric graph planes;
the spin structures come from the base frame.  The distinct slopes are
totally ordered by positive-definite difference.  We include a generating
set of scalar slopes and may enlarge the list as described below.
This order on slopes is separate from any directed ordering of their
Hamiltonian copies.

Denote the mirror images of $\Sigma_\ell$ and $G_M$ by $Q_\ell$ and $N_M$.
The graph images are line bundles after one common shift.  Put
$K_{\ell,M}=\Sigma_\ell\times G_M$.  For perfect complexes we use the Euler
pairing
$\chi(A,F)=\sum_j(-1)^j\dim\Ext^j(A,F)$.

\begin{proposition}[Realization with a self-extension bound]
\label{prop:realization}
Fix the spaces and test systems above.  Let $a\in H^n(D,\Q)$ satisfy
$\omega\smile a=0$, and define
\[
 \mathcal Z=\{x\in H^2(D,\Q):x\smile a=0\}.
\]
Suppose that cup product induces an injective map
\begin{equation}
 H^1(D,\Q)\otimes\mathcal Z\longrightarrow H^3(D,\Q).
 \label{setup:wedge-injective}
\end{equation}
After a field extension and an abstract embedding into $\C$, there is
$\mathcal E\in\Perf(B\times E^g)$ such that
\begin{equation}
 \dim\Ext^2(\mathcal E,\mathcal E)
 \le b_2(D)-\dim\mathcal Z.
 \label{setup:ext-bound}
\end{equation}
There is a common nonzero $\lambda\in\Q$, independent of $\ell,M$, with
\begin{equation}
 \chi(Q_\ell\boxtimes N_M,\mathcal E)
   =\lambda\,\langle a,[K_{\ell,M}]\rangle.
 \label{setup:euler}
\end{equation}
The mirror varieties, their enhanced categories, and the images of the
fixed test systems can be chosen before $a$ is specified.
\end{proposition}

The last assertion lets us choose divisor classes and ample classes on
the geometric mirror before choosing $a$.  One may use a preliminary
complex embedding to describe their lattices and cones: the geometric
N\'eron--Severi lattice, intersection pairing, and ampleness of its
classes are unchanged by algebraically closed field extension.  Later
embeddings therefore preserve the algebraic data used in the
construction.

The proof of Proposition~\ref{prop:realization} first constructs a graded
Lagrangian immersion representing a multiple of $a$.  Its degree-two
cohomology gives the right-hand side of \eqref{setup:ext-bound}.
Unobstructedness places it in a proper cyclic category.  We then identify
the subcategory generated by the embedded test objects with
$\Perf(B\times E^g)$ and project to that subcategory.  The categorical
completion is given in Section~\ref{sec:realization}.

\subsection{Sphere tests and the algebraic Mukai space}

In this subsection assume that $Y$ is a quartic K3 surface.
Seidel's quartic theorem gives the enhanced equivalence on a generating
category and split-generation by vanishing spheres
\cite[Theorem~1.3, Corollary~9.6, and Lemma~11.4]{SeidelQuartic}.
Rescaling the Novikov parameter accommodates the normalization of
$\omega_Y$.  Sphere branes are rational: no one-cycle holonomy condition
has to be imposed.  Compatibly graded isotopic spheres determine
quasi-isomorphic objects by \cite[Lemma~8.14]{SeidelQuartic}.  These facts
apply in particular to all the Hamiltonian copies used later.

We may add vanishing spheres from a generic pencil of quartics in
$\mathbf P^3$.  Symplectic transport identifies its smooth fibers with
$Y$, and the vanishing cycles span the kernel of ambient inclusion.
Their Poincar\'e duals therefore span $P_Y$.  Equivalently one can use
a hyperplane pencil in the corresponding Veronese embedding.

For any finite or countable set of the embedded sphere copies, there is
a common compatible almost complex structure with no nonconstant
holomorphic spheres and no nonconstant holomorphic disks on any
individual copy.  For somewhere injective curves the unmarked expected
dimensions are respectively $-2$ and $-1$.  Genericity, sphere
factorization, and the disk decomposition results of Kwon--Oh and
Lazzarini \cite{KwonOh,Lazzarini} give the assertion, as used in
\cite[Section~8]{SeidelQuartic}.  The torus graphs have vanishing relative
second homotopy, so a product almost complex structure has the same
exclusion property for the embedded product tests.

Write
\begin{equation}
 P_B=H^0(B,\Q)\oplus\NS(B)_{\Q}\oplus H^4(B,\Q),
 \label{setup:algebraic-mukai}
\end{equation}
with the Mukai pairing, in which $H^0$ and $H^4$ pair with a minus sign.
Before choosing a complex embedding, this denotes the geometric divisor
space together with its degree-zero and degree-four generators.
For an object $Q$ its Mukai vector is
$v(Q)=\operatorname{ch}(Q)\sqrt{\operatorname{td}(B)}$.

\begin{lemma}
\label{setup:rho-lemma}
The space $P_B$ has dimension $21$.  Sending the Poincar\'e dual of a
sphere test to the Mukai vector of its mirror image induces a rational
isometry
\begin{equation}
 \rho:P_Y\xrightarrow{\sim}P_B.
 \label{setup:rho}
\end{equation}
\end{lemma}

\begin{proof}
The Euler pairing of two sphere branes is minus their intersection
pairing.  On their mirror objects it is minus the Mukai pairing.
Since the test classes span the nondegenerate, $21$-dimensional space
$P_Y$, their Mukai vectors span a space of dimension at least $21$.

The mirror is the geometric generic member of the one-parameter family
obtained by resolving the diagonal-group quotient in Seidel's theorem.
That family has nonconstant period.  Before taking the finite symplectic
quotient, this is visible in the quartic pencil generated by the product
of the four coordinates and their fourth powers: at the Fermat member
the infinitesimal period in the product direction is nonzero, because
$x_0x_1x_2x_3$ survives in the Jacobian ring.  The group fixes the
holomorphic two-form.  Pullback by the quotient correspondence over an
open base therefore transports this variation to the resolved family.
Its geometric generic Picard rank is at most $19$ by period theory.
The HMS parameter is transcendental, including after its formal
reparametrization, and geometric Picard rank is unchanged by
algebraically closed field extension.  Hence $\dim P_B\le21$.

We have equality.  The pairing equality now also proves that the map on
sphere classes is well-defined: a relation among their classes maps to
a vector orthogonal to a spanning set of the nondegenerate space $P_B$,
and thus to zero.  The resulting surjective map preserves the pairings,
which proves \eqref{setup:rho}.
\end{proof}

\subsection{Graph images and full faithfulness on the test systems}

We use Abouzaid's family Floer theorem in its finite-object form
\cite[Theorem~2.10 and Appendix~A]{AbouzaidFaithful}.  It gives a faithful
$A_\infty$ functor to analytic perfect complexes.  For the present test
systems a dimension comparison will give full faithfulness.

The fibration $T_g\to\R^g/\Z^g$ is nonsingular, and both its base and
total space have vanishing second homotopy.  A graph induces
$b\mapsto(b,Mb)$ on fundamental groups, so
$\pi_2(T_g,G_M)=0$.  The graphs and their Hamiltonian copies are
therefore tautologically unobstructed.  The global zero section and
the base frame trivialize the section and spin twists in family Floer
theory.  Let $\mathfrak q$ be the Tate parameter, of positive valuation.
The rigid analytic mirror has multiplicative coordinates
$z=(z_1,\ldots,z_g)$ and periods $z\mapsto\mathfrak q^b z$, and is
$(E^g)^{\mathrm{an}}$.

\begin{lemma}[Graph images]
\label{setup:graph-images}
For each finite test system, the image of $G_M$ is a line bundle up to
shift, with numerical type $\varepsilon M$ for one common sign
$\varepsilon\in\{1,-1\}$.  After fixing the variance convention, all
shifts can be made equal.  The family Floer functor is fully faithful
on the finite systems and their Hamiltonian copies.
\end{lemma}

\begin{proof}
A graph has one transverse intersection with each torus fiber.  Its
local family Floer complex is consequently a rank-one module in one
degree.  We recall the numerical calculation of
\cite[Lemma~9.1]{A8}, including its extension to non-diagonal slopes.
On the real cover of the base lift a perturbed graph as
$p=Mq+\nabla f(q)$, where $f$ is periodic.  Comparison paths along this
lift have zero character offset: the fixed lattice offset of a lifted
graph boundary cancels at its two ends, and a residual fiber loop could
not bound.  Scalar trivializations on the real cover therefore give
the period multiplier
\begin{equation}
 j_M(b,z)=c_bz^{-\varepsilon Mb},\qquad
 b\in\Z^g,\quad c_b\in\Lambda^*.
 \label{setup:graph-multiplier}
\end{equation}
Here $\Lambda$ is the Novikov coefficient field and a change of variance
by perfect duality changes all numerical signs together.

Every integral symmetric matrix is a sum of signed matrices $uu^t$,
with $u\in\Z^g$.  Products of theta lines pulled back under
$E^g\to E$, $(x_j)\mapsto\sum u_jx_j$, have precisely the prescribed
type $\varepsilon M$.  Dividing their multiplier into
\eqref{setup:graph-multiplier} leaves a scalar character, hence a line
in $\operatorname{Pic}^0(E^g)$.  Projective rigid GAGA algebraizes the
image.  For the enhanced comparison one can use the common-refinement
\v Cech models of \cite[Section~9]{A8}; their restriction and cup maps
give the same dg comparison in $g$ variables.

For distinct slopes with definite difference, Floer cohomology has
dimension $|\det(M-M')|$ and is concentrated in one degree.  The quotient
of the corresponding lines is ample or anti-ample, and its total Ext
dimension is the same determinant.  Faithfulness is thus fullness for
these pairs.  For copies of the same slope, self-Floer cohomology has
the exterior-algebra dimensions of a $g$-torus.  A nontrivial
$\operatorname{Pic}^0$ quotient would have zero cohomology, contradicting
faithfulness.  The underlying lines of the copies therefore agree, and
their self-Ext dimensions again give fullness.

Choose the arrow direction in which definite graph comparisons have
degree zero.  The numerical sign must make the corresponding line
quotients ample.  Otherwise every pair among three ordered scalar
slopes would require the same nonzero shift difference $g$, which
contradicts additivity of shift differences.  With the correct sign,
all scalar shifts agree.  Comparison with extreme scalar slopes makes
the shift of every other graph the same.  Grading transport and the
self-cohomology support $[0,g]$ give the same conclusion for copies.
This is the normalization argument of \cite[Lemma~9.2]{A8}.
\end{proof}

\subsection{Finite generation and the order of choices}

\begin{lemma}
\label{setup:graph-generation}
Let $H$ be the principal product theta line on $E^g$.  There are
integers $k_1<\cdots<k_{g+1}$ such that, for every choice
$P_j\in\operatorname{Pic}^0(E^g)$, the lines
$H^{-k_j}\otimes P_j$ split-generate $\Perf(E^g)$.
The assertion persists under algebraically closed field extension.
\end{lemma}

\begin{proof}
Choose a vector-bundle generator $V$ of $\Perf(E^g)$
\cite{BondalVanDenBergh}.  Starting with $K_0=V$, take evaluation
surjections
\[
 0\longrightarrow K_j\longrightarrow
 W_j\otimes(H^{-k_j}\otimes P_j)
 \longrightarrow K_{j-1}\longrightarrow0.
\]
Here $W_j$ is a finite-dimensional vector space.  The exponents can be
chosen successively and uniformly in all preceding $P_j$: perform each
step over the proper parameter space
$(\operatorname{Pic}^0 E^g)^j$, using its Poincar\'e lines and relative
Serre generation and vanishing.  The kernels remain vector bundles.
The length-$g+1$ sequence gives a triangle
\[
 K_{g+1}[g]\longrightarrow C_V\longrightarrow V
 \longrightarrow K_{g+1}[g+1],
\]
where $C_V$ is a complex of sums of the displayed lines.  Since
$\Ext^{g+1}(V,K_{g+1})=0$, the triangle splits and $V$ is a summand of
$C_V$.  All evaluation sequences and base-change properties were
constructed over the parameter spaces, so the argument persists after
field extension.  The original generator remains a generator after
extension, for example by the usual perfect-diagonal criterion.
This is the argument of \cite[Lemma~9.3]{A8} in dimension $g$.
\end{proof}

If the sign convention gives ample rather than anti-ample powers,
dualize the preceding construction.  Thus the scalar graph tests can
always be chosen to include generators.  Further tests may be added
while retaining definite order: the integral symmetric matrices larger
than a fixed matrix in positive-definite order are Zariski dense.
Consequently every finite-dimensional span obtained by polynomial
evaluation on symmetric matrices can be obtained by continuing the
ordered list.  Extreme scalar slopes can then be added for the shift
normalization.

We fix these finite systems and their mirror images before choosing
$a$.  Generation in both factors survives the later field extensions.
Finally, $E$ is non-CM after every complex embedding used here.  Its
$j$-invariant has negative Novikov valuation, so is transcendental over
$\Q$, whereas a CM elliptic curve has an algebraic $j$-invariant.
